\section{Problem Statement}

Bug reports are an important source for fault localization.
A high-quality report may directly pinpoint faulty classes or methods; however, in practice, bug reports are often noisy, misleading, or incomplete.
Table~\ref{tab:example} illustrates this challenge using two real-world Apache Druid bugs.

In \textbf{Druid-3174}, the stack trace explicitly mentions \texttt{WorkerTaskMonitor}, yet the actual race condition resides in \texttt{RemoteTaskRunner}, where \texttt{WorkerTaskMonitor} merely serves as a triggering entry point.
Accurate localization therefore requires reasoning beyond the reported class.
In \textbf{Druid-1360}, the bug report contains no explicit class or method names at all, making it particularly difficult for existing fault localization techniques to establish reliable anchors in the codebase.
Without symbolic references, both IR-based and LLM-based approaches struggle to map the report to relevant program locations.

\vspace{-0.5em}
\begin{table}[h] \small
\centering
\caption{Bug Reports from Apache Druid} \label{tab:example}
\vspace{-8pt}
\begin{tabularx}{\textwidth}{X}
\Xhline{1.0pt}
\textbf{Title:} Overlord assigns too many tasks to middle manager \textit{\#3174} \\\hline
\textbf{Description:} RTR has a guard against multiple things trying to launch on the same worker via \\
workersWithUnacknowledgedTask.putIfAbsent(immutableZkWorker.get().getWorker().getHost(), ...\\
INFO [WorkerTaskMonitor] io.druid.indexing.worker.WorkerTaskMonitor - Submitting runnable for task ... \\
INFO [WorkerTaskMonitor] io.druid.indexing.worker.WorkerTaskMonitor - Affirmative. Running task ... \\ \hline
\textbf{FlexFL:} 1. \textit{find\_class}("WorkerTaskMonitor") $\rightarrow$ 2. \textit{get\_methods\_of\_class}("WorkerTaskMonitor") $\rightarrow$ 3. \textit{get\_code\_snippet\_of\_method}("mainLoop()") $\rightarrow$ 4. \textit{find\_class}("ForkingTaskRunner") ...\\[5pt]
\textbf{Top 1}: io.druid.indexing.worker.WorkerTaskMonitor.mainLoop() \\
\textbf{Top 2}: io.druid.indexing.overlord.ForkingTaskRunner.run(Task) ... \\
% \textbf{Top 3}: io.druid.indexing.common.task.Task.getId() ...\\
\Xhline{1.0pt}
\textbf{Title:} Race condition in autoscaling terminates nodes that were just assigned tasks \textit{\#1360}\\ \hline
\textbf{Description:} It is possible for the indexing service auto-scaling to terminate nodes that have just been assigned a task, but that have not updated the task status to running yet.
... If you are unlucky and autoscaling checks right after the task has been assigned, but right before the worker announces the task status, and the worker is not running any other tasks, your tasks will get killed. ...\\ \hline
\textbf{FlexFL:} 1. \textit{find\_class}("autoscaling") $\rightarrow$ 2. \textit{get\_methods\_of\_class}("AutoScaler") $\rightarrow$ \\ 3. \textit{get\_code\_snippet\_of\_method}("terminate(List<String>)") ...\\
\textbf{Top 1}: io.druid.indexing.overlord.autoscaling.AutoScaler.terminate(List<String>) \\
\textbf{Top 2}: io.druid.indexing.overlord.autoscaling.AutoScaler.terminateWithIds(List<String>) ... \\
\Xhline{1.0pt}
\end{tabularx}
\end{table}
\vspace{-0.5em}


We further examine FlexFL, a representative state-of-the-art LLM-based fault localization system that relies solely on bug reports.
FlexFL prompts the LLM to iteratively issue function calls—first resolving class names, then enumerating methods, and finally ranking code snippets.
While effective when bug reports precisely align with program symbols, this workflow degrades significantly under incomplete or misleading reports.
For instance, in Druid-3174, FlexFL remains confined to \texttt{WorkerTaskMonitor}, failing to identify the true fault in \texttt{RemoteTaskRunner}.
In Druid-1360, lacking explicit class names, it incorrectly localizes the fault to \texttt{AutoScaler}, whereas the actual issue lies in \texttt{SimpleResourceManagementStrategy} and \texttt{RemoteTaskRunner}.

These cases highlight three fundamental \textbf{limitations} of bug-report–only LLM-based fault localization for concurrent programs:

\textit{(1) Knowledge insufficiency.}
Bug reports frequently omit critical structural and behavioral information required for grounding LLM reasoning.
When program entities are missing or underspecified, the LLM is forced to infer context from incomplete natural language, leading to unstable or hallucinatory reasoning.
We observe that FlexFL may produce inconsistent localization results across multiple runs on the same report, issuing different function calls and retrieving different candidate methods.

\textit{(2) Ineffective search-space narrowing.}
When bug reports lack explicit program identifiers, LLM-based approaches have no principled way to progressively narrow the search space.
The LLM is therefore forced to issue broad or heuristic queries, often retrieving code only loosely related to the actual fault.
This describes \textbf{Druid-1360}, where the report contains no class or method names and FlexFL incorrectly focuses on \texttt{AutoScaler}, a component only tangentially related to the failure.
Without structured guidance or domain constraints, the LLM fails to identify the relevant concurrency context and repeatedly explores incorrect regions.

\textit{(3) Difficulty reasoning over deep and cross-thread call chains.}
Concurrent Java programs commonly exhibit deep, polymorphic, and cross-thread call hierarchies.
Due to limited context capacity, LLM-based approaches struggle to reason beyond local or downward call expansion.
In \textbf{Druid-3174}, although \texttt{WorkerTaskMonitor} appears in the report, the actual race occurs several layers downstream in \texttt{RemoteTaskRunner}.
Tracing such cross-thread interactions exceeds the practical reasoning depth of current LLM-based pipelines, causing them to remain trapped near the reported entry point.

\noindent
\textbf{Our Solutions.}
To overcome the above limitations, we propose three complementary design choices.

\textit{(1) Structured concurrency knowledge.}
To address knowledge insufficiency, we construct a \emph{Concurrent Knowledge Base (CKB)} that explicitly captures concurrency-relevant entities and behaviors, including thread entry points, shared variables, and synchronization mechanisms.
By grounding LLM reasoning in structured and statically extracted concurrency knowledge, the CKB reduces hallucination and stabilizes reasoning when bug reports lack precise program details.

\textit{(2) Hierarchical search-space narrowing.}
To enable effective localization under incomplete or misleading bug reports, we design a \emph{hierarchical retrieval strategy} that progressively narrows the search space from components to packages and then to thread entry points.
This coarse-to-fine localization process leverages the high cohesion and low coupling of software systems, allowing the LLM to focus on increasingly relevant concurrency contexts while keeping the exposed context bounded.

\textit{(3) Interaction-level reasoning without deep call-chain traversal.}
To overcome the difficulty of reasoning over deep and polymorphic Java call chains, we introduce an \emph{interaction-level DSL} that summarizes concurrency behavior in terms of shared-resource accesses across thread entries.
This representation allows the LLM to reason directly about potential interleavings and synchronization patterns, without explicitly traversing long call chains or enumerating execution paths.